\clearpage
\ExposeTOC
\begin{center}
{\Large\bfseries 报告 XXI}\par
\vspace{1.5cm}
{\Large\bfseries 完全交的上同调之层次}\par
\vspace{0.8cm}
{\large N. Katz}\par
\vspace{0.5cm}
附 P. Deligne 的附录 (\S 5)
\end{center}

\vspace{1.5cm}
\section*{目录}
\begin{description}
\item[0.] 引言 \hfill 1
\item[1.] 完全交的 Hasse--Witt 矩阵 \hfill 2
\item[2.] 有限定域上的余层次 \hfill 14
\item[3.] 对完全交的应用 \hfill 16
\item[4.] 结论 \hfill 19
\item[5.] 整性定理 \hfill 21
\end{description}

\clearpage
\section*{0. 引言}

设 $X$ 是维数 $n\geq3$ 的光滑射影簇, $Y(d)$ 是次数为 $d$ 的超曲面截面,
$E^{n-1}(Y(d),\mathbb Q_\ell)$ 是 $Y(d)$ 的“消灭”上同调
(XVIII~5.4.4). 本报告旨在证明: 当 $d$ 充分大且 $Y(d)$ 充分一般时,
$E^{n-1}(Y(d),\mathbb Q_\ell)$ 的全体并不都能通过代数对应“来自”光滑固有代数簇的
$H^i$, $i<n-1$. 在特征零的情形, 这立即由复上同调上的 Hodge 双分次以及凝聚上同调中的下述事实得出:
\begin{equation*}
\dim H^{n-1}(Y(d),\mathcal O_{Y(d)})\longrightarrow\infty
\qquad\text{当 }d\longrightarrow\infty.
\end{equation*}
印本中有明显的排印误文 \emph{et de démontrer} 和 \emph{suffisament};
上文采用的读法为 \emph{est de démontrer} 和 \emph{suffisamment}.

在特征 $p$ 的情形, 此方法同样把凝聚上同调群
$H^{n-1}(Y(d),\mathcal O_{Y(d)})$ 与平展上同调群
$H^{n-1}_{\mathrm{et}}(Y(d),\mathbb Q_\ell)$ 联系起来. 为此, 在有限基域上工作.
此时二者由 $Y(d)$ 的 zeta 函数联系起来: 该函数可用 $\ell$ 进上同调表示,
而它“模 $p$ 的约化”则用凝聚上同调表示 (参见 XXII~3.9).
对于完全交, 由于已详细了解这两种上同调 (平展上同调与凝聚上同调),
此方法的实施没有困难. 它把主要结果定理~4.2 归结为关于 Hasse--Witt 运算的结果,
即定理~1.4 (见 1.0).

A. Grothendieck 给出了基于晶体上同调的一般思想的启发式论证,
表明实质上用同一方法, 上述一般断言对每个光滑射影 $X$ 都应成立.

即使在完全交这一特殊情形, 证明也在 2.2 中实质地用到 Frobenius 特征值的整性定理~2.1;
该定理最近由 P. Deligne 证明, 这里在附录 (\S 5) 中给出证明.
印本称其为“附录 I”, 而标题和文献结构均表明 Deligne 的附录是 \S 5.

\section*{1. 完全交的 Hasse--Witt 矩阵}

\subsubsection*{定义 1.0.}
特征指数为 $p$ 的域 $k$ 上的固有概形称为\emph{特殊的 (special)}, 如果加性 $p$ 线性运算
\begin{equation}\tag{1.0.1}
\mathfrak F_p:\mathcal O_X\longrightarrow\mathcal O_X,
\qquad g\longmapsto g^p,
\end{equation}
诱导幂零运算
\begin{equation}\tag{1.0.2}
\mathfrak F_p:H^{\dim X}(X,\mathcal O_X)
\longrightarrow H^{\dim X}(X,\mathcal O_X).
\end{equation}
等价地, 若 $M=\dim_k H^{\dim X}(X,\mathcal O_X)$, 则 $p^M$ 线性运算
\begin{equation}\tag{1.0.3}
\mathfrak F_p^M=\mathfrak F_{p^M}:\mathcal O_X\longrightarrow\mathcal O_X,
\qquad g\longmapsto g^{p^M},
\end{equation}
将 $H^{\dim X}(X,\mathcal O_X)$ 化为零.

回顾 \textup{[2]}, 对每个 $i\in\mathbb Z$, 由 (1.0.1) 诱导的同态
\begin{equation*}
H^i(X,\mathcal O_X)\xrightarrow{\ \mathfrak F_p\ }H^i(X,\mathcal O_X)
\end{equation*}
通常称为 $H^i(X,\mathcal O_X)$ 上的 \emph{Hasse--Witt 同态}.
注意, 当 $p=1$ 时, $X$ 非特殊恰指 $H^{\dim X}(X,\mathcal O_X)\ne0$.
